\chapter{Discussion and Perspectives}
\label{chap:discussion}

{ \Large \leftwatermark{
		\put(-67,-66.5){ 1 }
		\put(-67,-91.5){ 2 }
		\put(-67,-116.5){ 3 }
		\put(-67,-141.5){ 4 }
		\put(-67,-166.5){ 5 }
		\put(-76.5,-199){\includegraphics[scale=0.8]{img/thumbindex/thumbindex_wongOrange.eps}} \put(-67,-191.5){ {\color{white} 6 }}
	} \rightwatermark{
		\put(350.5,-66.5){ 1 }
		\put(350.5,-91.5){ 2 }
		\put(350.5,-116.5){ 3 }
		\put(350.5,-141.5){ 4 }
		\put(350.5,-166.5){ 5 }
		\put(346.5,-199){\includegraphics[scale=0.8]{img/thumbindex/thumbindex_wongOrange.eps}} \put(350.5,-191.5){ {\color{white} 6 }}
}}

\newpage

\section{Summary of this thesis}
It is widely acknowledged that the ultimate goal of most (human) genetics research is to improve outcomes in healthcare. This can be done in many different ways, for example through drug discovery, biomarker identification, or better diagnostics, but also by improving our understanding of health and disease\cite{mcguireRoadAheadGenetics2020, claussnitzerBriefHistoryHuman2020e}. This thesis focuses on the latter, improving our understanding of health and disease, by mapping the genetic regulation of gene expression. These genotype-expression maps help with understanding the molecular mechanism with which genetic variants may contribute to disease risk.  

The regulatory networks, in which genes form pathways, are incredibly complex containing many layers of regulation. The function of a specific gene or the gene itself can both be candidates for drug targets. To identify these targets it is important to know how the genes function within complicated pathways. One way to approach this is by identifying which genetic factors affect the genes through eQTL studies. Subsequently building gene correlation networks to identify [AGM2]which pathways these genes are involved in. Preferably this should be done in cell types or tissues that are relevant for the phenotype of interest. In this thesis we have outlined the genetic architecture of molecular traits in \textbf{Chapter 1}\cite{claringbouldGeneticArchitectureMolecular2017} and identified the effect of rare variants on gene expression data in whole blood in \textbf{Chapter 2}\cite{kleinImbalancedExpressionPredicted2020}. We have also presented statistical methods that we developed to look at cell type interacting genetic effects in blood in \textbf{Chapter 4}\cite{raulaguirre-gamboaDeconvolutionBulkBlood2020}, and we have analysed the effects of common variants on gene expression in the brain in \textbf{Chapter 5}. 

This thesis chapter will outline how this data and these methods can be used to help prioritize candidate causal genes, assist in assessing the feasibility of drug compounds and create hypotheses for the role of identified genes on disease phenotypes, and how future work can improve on this.

\section{Understanding genetic regulation of gene expression}
Genome-wide association studies (GWAS) have allowed to gain valuable insights into the genomic loci that convey disease risk.  Thanks to the sharing of genetic information and GWAS summary statistics there are now many hundreds of datasets with disease associated variants available\cite{visscher10YearsGWAS2017}. However, the mechanism by which these variants heighten risk of a disease, for example the genes that are affected by such a variant, are often not known. 

\subsection{Understanding the effects of rare variants on gene expression}
As discussed in \textbf{chapter 2}, the effect sizes of GWAS variants become stronger for rare variants. Because eQTL and meQTL studies are of substantially lower sample size than GWAS studies, it is difficult to assess the effects of the rare variants on expression and methylation. Generally, these studies are done on common variants (minor allele frequency (MAF) $>$ 0.05), as sample sizes are smaller than those for GWAS studies due to higher sequencing costs. Even so, in yeast it has been shown that rare variants have a stronger effect on gene expression than common variants\cite{bloomRareVariantsContribute2019}. A study in humans of 472 genes showed that both very high and very low expressed genes had a high burden of rare variants\cite{zhaoBurdenRareVariants2016}. Evidently, identifying the effects of rare variants on gene expression is important for understanding the genetic mechanism of disease. Others have shown allele specific expression (ASE) effects in many different tissues<cite gtex> and that it can help with diagnosing<cite the muscle ase paper>, but generally these are done with limited number of samples or limited number of genes. In \textbf{chapter 3}, we performed allele specific expression in a BIOS, a large Dutch healthy population cohort with RNA-seq measurements from blood. Typically, when interested inrare variants, genotypes from whole genome sequencing (WGS) are preferable over chip-array genotypes. Since chip-array typically contain only common variants, and imputing rare variants is difficult<citation?>. Unfortunately, WGS data was only available for a subset of the BIOS samples. To still be able to assess the effects of rare variants on gene expression, we used the RNA-seq data to do genotype calling as was done previously by Deelen \emph{et al.}<citation>. Additionally, we phased the RNA-seq based genotypes, something not done with RNA-seq based genotypes before. Phasing SNPs gives information of which alleles are on the same haplotypes, which can improve allele specific expression measurements as counts from multiple SNPs can be combined. We used the WGS data of the subset of samples to validate both the genotypes and the phasing. Unfortunately, due to the fact that SNPs called with RNA-seq are only found in regions with expressed exons, the phased haplotypes were often small. Because of this, we could get a better estimate of the allelic expression, as we could combine counts from the same haplotype, but we were unable to phase variants outside of expressed genes and therefore could not assess the effects of regulatory genes on ASE.

Using this ASE approach we were able to assess the effects of rare variants on gene expression. We found 5 individuals with predicted high impact heterozygous variants in known autosomal dominant disease genes. It was surprising to find these variants, as the participants had no known genetic disorders. When assessing the allelic imbalance of these SNPs, we found that the reference allele was always much higher <fill in minimum ratio> expressed than the alternative allele, and thanks to that the gene expression was average of above average. One hypothesis of how this could happen is due to a positive feedback loop. Let's say there is a gene A and gene B. The expression of gene A is regulated by gene B, and higher expression of gene A causes gene B to down regulate gene A. In normal circumstances, both alleles of gene A produce a functional transcript and the expression of both alleles is about equal. However, if one of the alleles of gene A has e.g. a stop-codon variant, the transcript from that allele is degraded due to nonsense-mediated decay. Now, there is less gene A product, and therefore gene B stops down regulating gene A. Gene A starts producing more transcripts, but due to the nonsense mediated decay, only one allele produces a transcriptional product. When gene A and gene B are in homeostasis, one of the alleles is much higher expressed than the other. This is in line with the fact that we see an enrichment in ASE for genes containing stop-coding variants.

\subsection{Digging into cellular context}
To find causal links and to understand the variant mechanism, it is imperative to study effects in the right tissue and cell type. This is one of the main drivers behind developing Decon2 in \textbf{chapter 4}, which allows us to deconvolute cell type interacting eQTLs from bulk RNA-seq samples. However, there are some limitations to this approach. Firstly, although we compare favourably against other cell count prediction tools xCell\cite{aranXCellDigitallyPortraying2017} and CIBERSORT\cite{newmanRobustEnumerationCell2015}, there is still an error in the predicted cell count proportions. This error propagates when deconvoluting the cell type mediated eQTL effect from bulk tissue samples. Secondly, because we use cell count proportions there are significant correlations between cell types. These correlations makes it difficult to pinpoint the cell type specific effect. For example, if neutrophils and monocytes are anti-correlated to each other, and if an interaction effect is found in a neutrophils, it could also be that the effect is actually due to an opposite effect in monocytes. Thirdly, it is difficult to deconvolute the effect in sub-cell types, and for Decon2 we deconvoluted only the 6 main immune cell types. However, there are many more sub cell types that might be very relevant for understanding disease.

One method to overcome these three problems is to use single-cell RNA-seq\cite{tangMRNASeqWholetranscriptomeAnalysis2009}. Recently, this technique has been shown to be able to identify cell-type specific eQTLs\cite{vanderwijstSinglecellRNASequencing2018b}. Cell types can be assigned to single cells on the basis of expression levels of marker genes, removing the error of predicting cell counts. Because eQTLs are calculated within only those cells from a specific cell type, there is no problem of correlated cell count proportions. Thirdly eQTLs can be assessed even in rare cell types (although enough cells need to be available to have enough power to detect eQTLs).

\subsection{Disease dependent effects}
Genetic effects are not only dependent on the cellular context, but physiological context as well. Some eQTLs are dependent 

\subsection{Functional validation}
This work has been ... crispr

\subsection{The need for more diversity in sampling}
One major obstacle of transitioning to precision medicine is that most of the large genetic cohorts have been done in European (52\%) or Asian (21\%) populations\cite{sirugoMissingDiversityHuman2019}. This is mainly problematic for polygenic disorders where there the causative variants and tagging SNPs are in weak LD\cite{sirugoMissingDiversityHuman2019}. For example, a multi-ethnic meta-analysis for pulmonary function identified over 50 new loci by including individuals from African, Asian, and Hispanic/Latino ethnicities\cite{wyssMultiethnicMetaanalysisIdentifies2018}.  Additionally, certain SNPs may only be present in specific populations\cite{adalsteinsdottirberglindNationwideStudyHypertrophic2014}.  In \textbf{chapter 5}, enough Cortex samples from individuals from African and East Asian descent are included for eQTL analysis, and we find strong concordance of directionality of eQTL effect. Together this suggests that generally, the genetic effect on gene expression will be the same, but the frequency of alleles is different. 
[AGM18]
\subsection{Drug predictions}
What is not addressed in this thesis is the prioritisation of drug compounds using genomics information. One of the tools that can be used is [AGM19]the connectivity map\cite{lambConnectivityMapUsing2006}. It has measured gene expression levels in cell cultures that were treated by different drug compounds, and compared that to gene expression levels of cell cultures that were not treated by any drug compounds. [AGM20]Fisetin\cite{gibbsHumanGenomeProject2020}.

\subsection{Outside of gene expression}
This thesis has focussed on the genetic regulation of gene expression. 

\section{Future work}
Looking back, the transformation that the genomics and genetics fields have made over the last 20 years is amazing. There have been tremendous advances that have led to a much better understanding of the genetic basis of various diseases. These insights are now often driving translational research. For example, for the drug development pipeline, the number drug mechanisms with direct genetic support increases from 2\% at the preclinical stage to 8.2\% for approved drugs\cite{nelsonSupportHumanGenetic2015d}. Since for many diseases we know which variants are associated with increased risk thanks to the many large scale GWAS studies that have been done (currently the GWAS catalog\cite{macarthurNewNHGRIEBICatalog2017a} contains 4118 studies for 3418 traits), this information can be leveraged to improve drug development. However, to be able to do this in the most optimal way we need to understand the mechanism by which these variants drive disease risk, and to link them to causal genes and proteins.
With continued efforts in building comprehensive maps of genotype-phenotype relationships under various circumstances, further advancements in high throughput functional variant screenings such as with crispr-cas, and continued application of new insights from these efforts for new therapeutic options, we will see genetics play more and more important roles in healthcare and drug development. [AGM24]

Altogether, future studies aiming to improve our understanding of molecular effects of genetics need to combine a lot of different aspects. Ideally, they should be done with multi-ethnic individuals. If interested in a specific disease, this can be limited to populations with high prevalence. They should be done using single-cell technologies, as the cellular context of genetic effects are very important. With large<cite eqtlgen> and multi-tissue<cite gtex> studies having been completed with bulk data, findings from the single-cell studies can be replicated in bulk tissues if needed. Since physiology also plays an important role, the most informative studies would include both cases and healthy individuals. In case of e.g. an immune disease, stimulations should also be included. Finally, although this thesis focuses on genetic effects on gene expression, gene expression alone will not be able to capture the full cellular response. Although gene expression and protein levels are correlated <cite gene x protein study>, there is correlation in-between, and SNPs might have effect on chromatin state or methylation but not on gene expression directly. [AGM27]Therefore, measurements should include not only RNA-seq, but also proteomics, ATAC-seq or DNASEseq. Consequently, any result found should [AGM28]be functionally validated, for example using crispr-cas screening.


\subsection{Collaboration and open data}
Finally, a lot of the work in this thesis, and all of \text{chapter 5}, is built on publicly available data. It has to be highlighted that without the academic community being so open to sharing their data much of the genetic work discussed here would not be possible. It is great to see that there continue to be large scale academic collaborations, such as GTEx, public-private collaborations, such as the UKBiobank<cite ukbiobank> and Fingenn<cite finngen>, as well as large disease specific collaborations like Psychencode and amp-ad<cite psychencode>, as without that my research would be impossible. Hopefully and ideally, collaborations like these will continue to happen as future research hangs on this. 


\bibliographystyle{naturemag}
\bibliography{chapters/chapter6-discussion/chapter6-discussion}
